List of Novel Discoveries

1. Residual Seed \(\delta\)
Discovery that seven iterations of the Arcan(\(\tau\)) angular bridge \(\psi\) produce the practical generator
  \[
  7\psi=\frac{\pi}{3}+\delta,\qquad\delta\approx0.0051676144\ \mathrm{rad}.
  \]
\(\delta\) is the generative seed of the entire radial-fold hierarchy: small enough to preserve near-six-fold symmetry, non-zero enough to keep the orbit emergent from the continuous transcendental family.

2. Dual Method
Winding method* extracts residual intensity and produces the *residual bulge**.
Radial-symmetry folding* organises the same seed into the *residual hexagon** (Order-6 + inner intersection hexagon).
Both structures are dual geometric expressions of one oriented residual seed.

3. 6.5-Winding Trace as Origin of Residual Intensity
Explicit geometric trace linking the maximal winding (\(w=3.25\), \(\theta_{\max}=6.5\pi\)) terminal ray \(\arctan(2\pi)\approx1.412965136507\) to the thinnest residual altitude \(1.490969476641\).
Linear gap
  \[
  1.490969476641-1.412965136507=0.078004340134
  \]
  (ratio \(\approx1.055206\)) is the first quantitative signature connecting the winding to the residual triangle.

4. Thinnest Equilateral-Equidistributal Transcendental Scaling Triangle
Identification of the residual triangle \(T\) (sides \(1.721623257385\) / \(1.737195272475\), altitude \(1.490969476641\), shoelace area \(1.2988990741\)) as the elementary residual cell.
Characterised simultaneously as:
  Thinnest (altitude compressed by \(\delta\)),
  Near-equilateral,
  Equidistributal (minimal cell of a dense residual orbit),
  Transcendental (metric data are transcendental invariants of the Arcan(\(\tau\)) family),
  Scaling (self-similar unit across the hierarchy).

5. Core Relations
Dimensionless residual intensity
  \[
  \rho=\frac{A_{\mathrm{slice}}}{\operatorname{Area}(T)}\approx0.472886
  \]
  derived from leftover area \(1.8426935795\) and residual slice \(0.6142311932\).
Area-derived residual bulge
  \[
  \operatorname{bulge}=c\cdot\rho\approx0.08512\qquad(c=0.18).
  \]
These Core Relations are invariant and propagate unchanged through every order.

6. Order-6 as Geometric Locus of the Core Relations
Order-6 is the unique locus at which the Core Relations first become visible simultaneously as:
  continuous edge curvature (residual bulge),
  discrete inner hexagonal measure (intersection hexagon of mean radius \(\approx0.57739\)).

7. Inner Hexagon from Intersecting Corner Points
Measurement of the secondary hexagon formed by the six intersection points of the interlocking residual triangles:
  mean side length \(\approx0.577330\),
  mean radius from origin \(\approx0.577390\).

8. Hierarchy as Pure Densification of a Single Residual Intensity
Demonstration that Orders 10 and 30 (and all higher densifications) introduce no new intensity.
They are pure combinatorial densifications of the identical residual intensity \(\rho\) fixed at Order-6.
The hierarchy contains no further elementary geometric invention—only progressive unfolding of one residual intensity.

9. Phase Windings as Directional Propagation of Topology
Topological formulation: phase windings are the oriented propagation of lines; geometry is the global organisation of those oriented lines.
The residual seed \(\delta\) is an oriented defect; residual bulge and residual hexagon are dual readings of one directional propagation.

10. Transcendental Invariance of Metric Data
Side lengths, altitude and shoelace area of \(T\) are transcendental invariants of the Arcan(\(\tau\)) family.
They descend from \(\arctan(\tau)\), \(\arctan(\pi)\) and \(\psi=\theta-\phi\) by operations that preserve transcendence.

11. Unified Cascade
Continuous Arcan(\(\tau\)) family  
→ seventh multiple near \(\pi/3\)  
→ residual seed \(\delta\)  
→ 6.5-winding altitude gap  
→ thinnest residual triangle \(T\)  
→ residual intensity \(\rho\) and bulge  
→ Order-6 radial fold + inner hexagon  
→ pure densification (Orders 10, 30, \ldots)  
→ outer \(E_8\) Petrie-ring geometry  
while remaining phase-locked to the Cosmological Clock.

These eleven results constitute the novel discoveries crystallised across the session. All of them descend from the single residual seed \(\delta\) generated by the continuous Arcan(\(\tau\)) family.
